\documentclass[11pt]{article}
\usepackage[utf8]{inputenc}
\usepackage[margin=1in]{geometry}
\usepackage{amsmath, amssymb, amsthm}
\usepackage{booktabs}
\usepackage{graphicx}
\usepackage{hyperref}
\usepackage{enumitem}
\usepackage{caption}
\usepackage{xcolor}

% Set link colors for professional appearance
\hypersetup{
    colorlinks=true,
    linkcolor=blue,
    filecolor=magenta,
    urlcolor=cyan,
    pdftitle={PCA of Global Markets - Ed Hayes},
}

% Path configuration to look one level up for the plots directory
\graphicspath{ {../plots/} }

\title{ML in Finance: Principal Component Analysis of Global Markets}
\author{Ed Hayes}
\date{March 2026}

\begin{document}

\maketitle

% --- CLICKABLE TABLE OF CONTENTS ---
\pagenumbering{roman}
\tableofcontents
\newpage
\pagenumbering{arabic}

\section{Introduction}
Principal Component Analysis (PCA) is a foundational dimensionality reduction technique used in quantitative finance to identify latent drivers of asset returns. Whether analyzing the term structure of interest rates or the sector rotation of the S\&P 500, PCA allows us to distill hundreds of correlated variables into a few interpretable factors.

\section{Mathematical Foundation}
\subsection{Data Representation and Stationarity}
To ensure the data is stationary, we transform the raw prices or yields:
\begin{itemize}
    \item \textbf{Interest Rates:} We utilize \textit{first differences} ($\Delta y_{t}=y_{t}-y_{t-1}$) because yields are already expressed in percentages.
    \item \textbf{Equities:} We utilize \textit{percentage changes} or returns ($r_{t}=\frac{P_{t}-P_{t-1}}{P_{t-1}}$) to normalize for absolute price levels.
\end{itemize}

\subsection{The Covariance Matrix}
PCA seeks to find the directions of maximum variance. Let $X$ be the $T \times n$ matrix of centered changes. The sample covariance matrix $\Sigma$ is:
\[ \Sigma=\frac{1}{T-1}X^{T}X \]

\subsection{Eigen-Decomposition}
We diagonalize $\Sigma$ to find the principal components:
\[ \Sigma=Q\Lambda Q^{T} \]
where $\Lambda$ contains the eigenvalues and $Q=[v_{1},v_{2},v_{3}]$ contains the eigenvectors.

\section{Cross-Asset Interpretation}
The first three principal components often explain over 95\% of market movement.

\begin{table}[h]
\centering
\begin{tabular}{@{}lll@{}}
\toprule
\textbf{Component} & \textbf{Fixed Income Interpretation} & \textbf{Equity Interpretation} \\ \midrule
PC1 & Level (Parallel Shift) & Market Factor (Beta) \\
PC2 & Slope (Twist) & Sector/Style Rotation \\
PC3 & Curvature (Butterfly) & Specific/Thematic Clusters \\ \bottomrule
\end{tabular}
\caption{Comparison of latent factors.}
\end{table}

\newpage
\section{Historical Case Studies}

\textcolor{red}{\textbf{These interpretations were created by Google Gemini. As such, you need to analyze them carefully and assess their accuracy for yourself.}}

\subsection{Evolution of the Interest Rate "Level" Factor (PC1)}
A longitudinal analysis of the Scree plots in Appendix \ref{app:ir_scree} reveals how the dominance of the Level factor (PC1) fluctuates with Fed policy and economic stability:
\begin{itemize}
    \item \textbf{2000--2005 (81.2\%):} High synchronization following the Dot-com bubble burst and subsequent easing cycle.
    \item \textbf{2005--2010 (60.9\%):} A significant drop in PC1 dominance. As seen in Appendix \ref{app:ir_scree}, the Slope factor (PC2) spiked to 26.3\% variance as the curve twisted violently during the Global Financial Crisis.
    \item \textbf{2010--2015 (86.6\%):} Maximum PC1 dominance. During the "Zero Bound" era, the entire curve was pinned by quantitative easing, leading to highly parallel movements.
    \item \textbf{2015--2020 (78.3\%):} A return to moderate levels as the Fed attempted to normalize rates before the COVID-19 shock.
    \item \textbf{2020--2025 (73.2\%):} While still dominant, high-dispersion events prevented the Level factor from reaching the levels seen in 2010.
\end{itemize}

\subsection{The 2008 Global Financial Crisis (2005--2010)}
During this era, the \textbf{Financials (XLF)} sector was the primary driver of market variance.
\begin{itemize}
    \item \textbf{Market Sensitivity:} XLF exhibited the highest positive weight in PC1 (Market Beta), as seen in the Factor Loadings in Appendix \ref{app:eq_loadings}.
    \item \textbf{Sector Specifics:} The massive spike in XLF within PC2 (Value vs. Growth) highlights that the crisis was specifically a financial-sector event.
\end{itemize}

\subsection{The Pandemic and Inflation Era (2020--2026)}
The current regime is defined by \textbf{Technology (XLK)} dominance and a clear Growth-led market structure.
\begin{itemize}
    \item \textbf{Tech as Beta:} In Appendix \ref{app:eq_loadings}, XLK shows extreme sensitivity in PC1, becoming a proxy for the overall market.
    \item \textbf{Growth vs. Value:} The extreme negative loading of XLK in PC2 illustrates its role as the engine for "Growth" trades.
\end{itemize}

\newpage
\appendix

% --- APPENDIX A: RATES SCREE ---
\section{Interest Rate: Scree Plots (2000--2025)} \label{app:ir_scree}
\begin{figure}[h!] \centering
    \includegraphics[width=0.65\textwidth]{Rates_2000_2005_Scree.png} \\ \vfill
    \includegraphics[width=0.65\textwidth]{Rates_2005_2010_Scree.png} \\ \vfill
    \includegraphics[width=0.65\textwidth]{Rates_2010_2015_Scree.png}
\end{figure} \newpage
\begin{figure}[h!] \centering
    \includegraphics[width=0.7\textwidth]{Rates_2015_2020_Scree.png} \\ \vspace{2cm}
    \includegraphics[width=0.7\textwidth]{Rates_2020_2025_Scree.png}
\end{figure} \newpage

% --- APPENDIX B: RATES LOADINGS ---
\section{Interest Rate: Factor Loadings (2000--2025)} \label{app:ir_loadings}
\begin{figure}[h!] \centering \includegraphics[width=0.9\textwidth]{Rates_2000_2005_Loadings.png} \end{figure} \newpage
\begin{figure}[h!] \centering \includegraphics[width=0.9\textwidth]{Rates_2005_2010_Loadings.png} \end{figure} \newpage
\begin{figure}[h!] \centering \includegraphics[width=0.9\textwidth]{Rates_2010_2015_Loadings.png} \end{figure} \newpage
\begin{figure}[h!] \centering \includegraphics[width=0.9\textwidth]{Rates_2015_2020_Loadings.png} \end{figure} \newpage
\begin{figure}[h!] \centering \includegraphics[width=0.9\textwidth]{Rates_2020_2025_Loadings.png} \end{figure} \newpage

% --- APPENDIX C: EQUITY SCREE ---
\section{Equity Sector: Scree Plots (2005--2025)} \label{app:eq_scree}
\begin{figure}[h!] \centering
    \includegraphics[width=0.65\textwidth]{Equity_2005_2010_Scree.png} \\ \vfill
    \includegraphics[width=0.65\textwidth]{Equity_2010_2015_Scree.png} \\ \vfill
    \includegraphics[width=0.65\textwidth]{Equity_2015_2020_Scree.png}
\end{figure} \newpage
\begin{figure}[h!] \centering
    \includegraphics[width=0.7\textwidth]{Equity_2020_2025_Scree.png}
\end{figure} \newpage

% --- APPENDIX D: EQUITY LOADINGS ---
\section{Equity Sector: Factor Loadings (2005--2025)} \label{app:eq_loadings}
\begin{figure}[h!] \centering \includegraphics[width=0.9\textwidth]{Equity_2005_2010_Loadings.png} \end{figure} \newpage
\begin{figure}[h!] \centering \includegraphics[width=0.9\textwidth]{Equity_2010_2015_Loadings.png} \end{figure} \newpage
\begin{figure}[h!] \centering \includegraphics[width=0.9\textwidth]{Equity_2015_2020_Loadings.png} \end{figure} \newpage
\begin{figure}[h!] \centering \includegraphics[width=0.9\textwidth]{Equity_2020_2025_Loadings.png} \end{figure}

\end{document}